\begin{textsample}{NEG Dim 4 – pi – Score -3.00 – pi/pi\_em\_11.txt}  \label{ex:f4_neg_242}
3.1.2 Ensino e aprendizagem O processo de ensino e aprendizagem tem sido objeto de reflexão há décadas.
A preocupação com o ato de educar vem desde a antiguidade.
O filósofo grego Sócrates, nos anos entre 469 a 399 a.C., opondo -se à filosofia que procurava explicar o mundo baseado na observação da natureza, defendeu a proposição de que o ser humano deveria conhecer a si mesmo, tendo como inquietação, levar as pessoas, por meio do autoconhecimento, à sabedoria e à prática do bem.
Para ele a valorização do diálogo era algo importantíssimo, pois serviria como meio de investigação acreditando, também, que a aprendizagem deveria começar a partir da educação do corpo, que por sua vez, permitiria controlar o físico.
Por volta dos anos 427 a 347 a.C., a formação do homem moral, como objetivo final da educação, assim como a defesa de um estado justo para acolher cada indivíduo foram objetivos defendidos por Platão.
Isso fez com que sua pedagogia se aproximasse da sua própria filosofia, dizendo que a busca da verdade é mais importante que a dos dogmas incontestáveis.
Este pensador defendeu a ideia de que professor ou aluno pensem sobre si próprios.
Aristóteles ( 384 a 322 a.C. ), discípulo de Platão, acreditava que a educação deveria ser tomada como o real caminho para a vida pública, tendo em vista que somente com uma adequada educação seria possível formar o caráter do aluno.
Neste sentido, este filósofo defendia a ideia de que a criança aprende com o ato de imitar os bons hábitos, e consequentemente, se tornar um cidadão ético e voltado à efetividade do bem comum em seu meio social.
No período medieval, com o paradigma teocêntrico em alta nos mais diversos segmentos ( religioso, cultural, político, epistemológico ), a educação era tratada como especificidade da Igreja, assim como a expansão da mesma na Europa.
Era responsabilidade da Igreja a disseminação da cultura, bem como dos processos educacionais deste contexto histórico, destacando -se um prévio curriculum instrucional embasado nas sete artes liberais, compostas por : gramática, retórica, lógica, aritmética, geografia, astronomia e música.
A Igreja torna -se a detentora do conhecimento e de formas educacionais deste período, tendo em vista que somente esta instituição poderia dar aval, conforme seus interesses, do que poderia ser ou não lecionado.
O ensino era pautado na leitura de textos e exposição das principais ideias pelo professor, formando muitas vezes debates entre professores e alunos ( scholastica disputattio ), e que se tornou conhecido como escolástica, método de estudos muito usado por São Tomás de Aquino ( 1225-1274 ).
No início da modernidade, é possível ser observada uma importante transição no que se refere aos métodos educacionais, onde o método escolástico ( medieval ) ganha uma remodelagem fundamentando -se não mais no conhecimento revelado ( religioso ), mas na vertente que defende que o próprio ser humano, enquanto ser racional, tem aparato mental e pode estabelecer critérios específicos para o alcance do conhecimento verdadeiro.
Este intermediado pelas experiências e uso coerente da razão, não mais tendo a razão como mero instrumento da religião.
Neste sentido, podem -se destacar algumas perspectivas educacionais que têm sua origem e motivações a partir de significativos acontecimentos e mudanças ocasionados em determinados momentos da história moderna, influenciando os mais diversos aspectos antropológicos, tais como econômico, político, social e cultural.
Com isto, pode ser afirmado que na modernidade têm -se ênfases de propostas educacionais direcionadas a métodos de ensino pensados a partir de fundamentos filosóficos, conforme se observa na continuidade deste texto.
Em uma perspectiva educacional pós contexto da Revolução Francesa – na qual se defendeu uma sociedade movida pelos princípios da Liberdade, Igualdade e Fraternidade, como direito a ser usufruído por toda sua população, surge uma marcante discussão a partir das ideias de Karl Marx ( 1818-1883 ), onde se tem a defesa de um ensino pautado na transformação social.
Na visão deste autor a sociedade pode ser melhor compreendida como uma composição de relações sociais, onde o ser humano é feito a partir de suas relações com a alteridade, ou seja, com os outros homens.
Assim, o ser humano seria movido à base de suas necessidades e liberdade, existindo uma dialética em que se possam justificar os fins.
No século XX, teve grande contribuição na formação da visão pedagógica, os pensamentos de Maria Montessori ( 1870-1952 ), ao defender a ideia de que a criança conquistava a educação, partindo da concepção de que já nascemos com a capacidade de ensinar a nós mesmos, desde que nos sejam dadas as condições.
Seus estudos basearam -se na observação de que as crianças aprendem melhor pela experiência direta da procura e descoberta, cabendo ao professor acompanhar o processo e detectar o modo particular de cada um manifestar seu potencial.
Ao longo do século XX, diferentes concepções pedagógicas fundamentaram o processo de ensino e aprendizagem, tais concepções, segundo Saviani ( 2006 ), podem ser agrupadas em duas grandes tendências : a primeira seria composta pelas concepções pedagógicas que dariam prioridade à teoria sobre a prática e a segunda tendência, inversamente, compõe -se das concepções que subordinam a teoria à prática e dissolvem a teoria na prática.
As modalidades de pedagogia tradicional, de vertente religiosa ou leiga, dão prioridade à teoria sobre a prática, e focam especificamente nas teorias de ensino.
Ao passo que as modalidades de pedagogias novas priorizam a prática, cuja ênfase é posta nas “ teorias da aprendizagem ”.
Ainda conforme Saviani ( 2006 ), as concepções tradicionais, desde a pedagogia de Platão e a pedagogia cristã, passando pelas pedagogias dos humanistas, assim como a pedagogia idealista de Kant, Fichte e Hegel, o humanismo racionalista, a teoria da evolução pautavam -se pela centralidade da formação intelectual.
Aqui seria possível observar que o protagonista do conhecimento era o professor, responsável por transmitir os conhecimentos acumulados pela humanidade segundo uma \textbf{gradação} lógica, cabendo aos alunos assimilar os conteúdos que lhes são transmitidos.
Nesse contexto, a prática era determinada pela teoria que a moldava, fornecendo -lhe tanto o conteúdo como a forma de transmissão pelo professor.
No que se refere às correntes educacionais ligadas à renovação metodológica, temos como referenciais pensadores como Rousseau, Kierkegaard, Stirner, Nietzsche e Bergson, chegando ao movimento da Escola Nova, bem como às pedagogias não diretivas, à pedagogia institucional e ao construtivismo priorizando em suas teorias da aprendizagem uma radical mudança, no qual a centralidade é o educando.
Este como protagonista da aprendizagem que, por meio da interação entre si e o professor, constroem seus conhecimentos.
Com esta defesa de um método voltado mais aos educandos, ao professor caberia o papel de acompanha -los de modo a auxiliá -los em seu próprio processo de aprendizagem.
Nessa perspectiva de aprendizagem, o primado é sobre a prática, o eixo do trabalho pedagógico desloca -se da compreensão intelectual para a atividade prática ; dos conteúdos cognitivos para os métodos ou processos de aprendizagem ; do professor para o aluno ; do esforço para o interesse ; da quantidade para a qualidade ( SAVIANI, 2006 ).
Essa tendência torna -se hegemônica sob a forma do movimento da Escola Nova, na segunda metade do século XX, e assume novas versões, entre as quais o construtivismo, sendo a mais difundida na atualidade.
A partir do século XX, a ênfase do processo de aprendizagem se desloca para os métodos, estabelecendo o primado dos fundamentos psicológicos da educação.
Nesse, o professor conduz a aprendizagem, mas o conhecimento tem como ponto de partida a experiência já existente ou a ser realizada pelo próprio aluno, conforme a perspectiva de Piaget ( 1983 ).
Deste modo, pode -se entender que a aprendizagem, no contexto teórico construtivista, está subordinada ao desenvolvimento, ou seja, é sempre provocada por uma situação e depende do desenvolvimento intelectual e da estrutura da própria inteligência.
Na perspectiva construtivista de Piaget, o conhecimento humano se constrói na interação homem-meio, sujeito-objeto.
O aprendizado depende essencialmente de uma prática pedagógica adequada, que favorecesse essa interação, de forma a que gradativamente o aluno pudesse internalizar esse conhecimento, isto é, assimilá -lo integrando -o aos seus demais conhecimentos ( PIAGET, 1975 ).
Conhecer consiste em operar sobre o real e transformá -lo a fim de compreendê -lo, é algo que se dá a partir da ação do sujeito sobre o objeto de conhecimento.
As formas de conhecer são construídas nas trocas com os objetos, tendo uma melhor organização em momentos sucessivos de adaptação ao objeto.
Essa adaptação ocorre através da organização, sendo que o organismo discrimina entre estímulos e sensações, selecionando aqueles que irá organizar em alguma forma de estrutura.
A adaptação possui dois mecanismos opostos, mas complementares, que garantem o processo de desenvolvimento : a assimilação e a acomodação.
Segundo Piaget, o conhecimento é a equilibração/reequilibração entre assimilação e acomodação, ou seja, entre os indivíduos e os objetos do mundo.
A aprendizagem, conforme defende Vygotsky ( 2001, p. 115 ), “ [ … ] pressupõe uma natureza social específica e um processo através do qual as crianças penetram na vida intelectual daqueles que a cercam ”.
É numa relação intrínseca do sujeito com o meio físico e social, orientada por instrumentos e signos ( entre eles a linguagem ), que se processa o seu desenvolvimento cognitivo.
Nessa perspectiva, o desenvolvimento psíquico do homem se realiza por meio de processo de internalização ( VYGOTSKY, 2001 ).
Segundo esse autor, as relações intrapsíquicas ( atividade individual ) constituem -se a partir das relações interpsíquicas ( atividade coletiva ).
É neste movimento do social ao individual que se dá a apropriação de conceitos e significações, ou seja, dá -se a apropriação da experiência social da humanidade.
Desta forma, podemos entender que a aprendizagem não ocorre espontaneamente, ela é mediada culturalmente.
Nas palavras de Leontiev ( 1978, p. 264 ) “ O homem não está evidentemente subtraído ao campo da ação das leis biológicas.
O que é verdade é que as modificações biológicas hereditárias não determinam o desenvolvimento sócio-histórico do homem e da humanidade ”.
A aprendizagem se dá por meio da interação e construção do conhecimento e também dos laços afetivos entre pessoas.
Wallon argumenta que as trocas relacionais da criança com os outros são fundamentais para o desenvolvimento da pessoa.
O meio é um \textbf{complemento} indispensável ao ser vivo.
Ele deverá corresponder às suas necessidades e as suas aptidões sensório-motoras, depois \textbf{psicomotoras}.
Não é mesmo verdadeiro que a sociedade coloca o homem em presença de novos meios, novas necessidades e novos recursos que aumentam a possibilidade de evolução e diferenciação individual.
A constituição biológica da criança ao nascer não será a única do seu destino ( … ).
Os meios em que vive a criança e aqueles com que ela sonha constituem a forma que amolda sua pessoa ( … ) ( WALLON, 1975, p. 164 ).
Para Emília \textbf{Ferreiro} ( 1983 ), a construção do conhecimento tem uma lógica individual na escola e fora dela.
A aprendizagem não é provocada pela escola, mas pela própria mente das crianças, em que esse processo deve ser \textbf{gradual} dependendo de sua assimilação e de uma reacomodação dos esquemas internos, o que demanda tempo.
O processo inicial da aprendizagem é explicado também por variáveis sociais, culturais, políticas e psicolinguísticas.
Complementa a autora, que não existe ponto zero da aprendizagem escrita ; a criança sempre apresenta um conhecimento prévio que o sujeito reestrutura a partir de um processo de acomodação e assimilação mental.
Segundo Antoni Zabala ( 1998 ), a finalidade da escola é promover a formação integral dos alunos.
Para ele, é na instituição escolar, através das relações construídas a partir das experiências vividas, que se estabelecem os vínculos e as condições que definem as concepções pessoais sobre si e os demais.
Daí a necessidade de uma reflexão profunda e permanente da condição de cidadania dos alunos, e da sociedade em que vivem.
Zabala atenta ainda para as particularidades dos processos de aprendizagem de cada aluno.
Na visão de Libâneo ( 1994, p. 88 ), > [ … ] aprender é o processo de assimilação de qualquer forma de conhecimento, daquilo bem simples onde a criança aprende a manipular os brinquedos, aprende a fazer contas, lidar com as coisas, nadar, andar de bicicleta etc., até aquilo mais complexo onde uma pessoa aprende a escolher uma profissão, lidar com as outras.
Dessa forma as pessoas estão sempre aprendendo.
Na origem da aprendizagem está todo um processo de assimilação no qual o aluno, com a orientação do professor, passa a compreender, refletir e aplicar os conhecimentos adquiridos, assim a aprendizagem é observada com a colocação em prática, por parte do estudante, desses conhecimentos que foram transmitidos durante uma aula ou atividade.
A aprendizagem é algo que modifica o pensamento.
Para que se possa haver a aprendizagem o aluno necessita ser estimulado com conteúdo de seu alcance, textos que tratem de sua realidade.
Nas palavras de Libâneo ( 1994, p. 90 ) “ a relação entre ensino e aprendizagem não é algo mecanizado, não é apenas uma transmissão feita pelo professor que ensina para um aluno que aprende ”.
Tem para ele, ao contrário, um dever de “ estabelecer exigências e expectativas que os alunos possam cumprir e, com isso, mobilizem suas energias.
Esse sim, pois é o seu papel, impulsionar a aprendizagem, precedendo -a muitas das vezes. ” Segundo Paulo Freire ( 1996, p. 14 ), não existe ensino sem aprendizagem, e nesse processo, “ educador e educandos, lado a lado, vão se transformando em reais sujeitos da (re)construção do saber, pois o conhecimento não está no professor, o conhecimento circula, é compartilhado ”.
O processo de ensinar tem por finalidade o alcance da aprendizagem pelo estudante.
Nesse sentido, na condução satisfatória do processo de ensino, efetivo e que agregue valor a vida dos estudantes, o professor adquire o papel de mediador e facilitador do processo de ensino-aprendizagem, assim precisa utilizar métodos adequados e centrados na realidade dos estudantes, num diálogo constante entre estes atores, buscando estabelecer e fortalecer vínculos.
Esse estreitamento de vínculos deve ser ponderado por regras onde todos têm liberdade, uma vez que, como afirma Freire ( 1986. p.54 ) : “ uma experiência dialógica que não se baseia na seriedade e na competência é muito pior do que uma experiência ‘ bancária ’, onde o professor simplesmente transfere conhecimento ”.
Com isto, pode ser afirmado que nesse diálogo, o encontro das pessoas ocorre para que elas reflitam sobre o mundo, sobre como transformá -lo e melhorá -lo, tendo na sala de aula o seu primeiro lugar de exercício prático.
Frente a isto, objetiva -se uma educação que ultrapasse a mera formação conceitual, conduzindo o alunado à reflexão sobre seu papel na sociedade, para que ele possa nessa relação teoria e prática encontrar novas maneiras de atuar no contexto em que vive.
Uma Educação que tenha a intenção de formar no aluno uma atitude crítica quanto às relações sociais nas quais se insere.
Constatar essas relações possibilita colocar “ nas mãos dos educadores uma arma de luta capaz de permitir -lhes o exercício de um poder real, ainda que limitado ” ( SAVIANI, 2003, p.31 ), uma vez que a educação se torna decisiva neste constante processo de transformação da realidade social.
A criação de sentido pessoal e social, orientada por princípios de solidariedade e favorecida pelas interações que aumentam o aprendizado instrumental produz a aprendizagem dialógica.
Essa aprendizagem apresenta -se como ferramenta que vai auxiliar na superação de desafios que estão internalizados ao bom êxito do processo de ensino e aprendizagem, visto que, essa aprendizagem fundamenta -se e expressa -se por meio do diálogo igualitário e na interação mais diversa possível entre as pessoas.
Como na contemporaneidade muito se tem refletido a respeito de como deve ser a atuação do professor em sala de aula no processo de ensino/aprendizagem, reafirmamos que cabe ao professor, na coordenação do processo de aprendizagem, articular os conhecimentos científicos, tecnológicos e artísticos, cientes dos vários elementos didáticos que influenciam e possibilitam a apropriação dos conhecimentos pelos alunos, seguindo uma perspectiva dialética entre o saber imediato e o mediato.
Com adoção de prática pedagógica que tenha dinâmica própria, que permita o exercício do pensamento reflexivo, conduza a uma visão política de cidadania e que seja capaz de integrar a arte, a cultura, os valores e a interação, propiciando, assim, a autonomia dos alunos e a reflexão sobre seu lugar no mundo, de forma significativa.

% matched lemmas: complemento, ferreiro, gradação, gradual, psicomotor
\end{textsample}
